\documentclass[../../../include/open-logic-section]{subfiles}

\begin{document}

\olfileid{sth}{z}{arbintersections}
\olsection{ضميمه: بندښت، ادراک او تقاطع}

په \olref[sth][z][milestone]{sec} کښې مو وړانديز وکړ چې
د \olref[sfr][][]{part} ساده کار ته بېرته وګورئ او ويې څېړئ
چې په~$\Zminus$ کښې ترسره کېدلے شي. که مو دا مشوره تعقيب کړې وي،
ښايي \emph{يوه} خبره مو ستونزمنه ليدلې وي: د ډېډېکېنډ د
\emph{بندښتونو} په بحث کښې د \emph{تقاطع} کارول.

د \olref[sfr][infinite][dedekind]{Closure} له مخې ياد ولرئ چې
\begin{align*}
  \closureofunder{f}{o} & = \bigcap\Setabs{X}{o \in X \text{ او $X$ د
$f$ له مخې بند دے}}.
\intertext{د دې عمومي بڼه د لاندې ډول يو تعريف دے:}
  C & = \bigcap\Setabs{X}{\phi(X)}.
\end{align*}
خو دا د خبرداري زنګ دے: ځکه چې ساده ادراک ناکام دے، نو هېڅ
تضمين نۀ شته چې $\Setabs{X}{\phi(X)}$ شتون لري. په دې توګه
داسې تعريفونه په خطرناک ډول د \emph{چل} په شان ښکاري.

خو له نېکه مرغه دا چل نۀ دے؛ يا که په اوسنۍ بڼه \emph{چل}
وي، نو په صادقانه کار سره ئې سم جوړولے شو. د دې کار نښه په
\olref[sth][z][sep]{prop:intersectionsexist} کښې وړاندې شوې وه،
چې ولې $\bigcap A$ د هر $A \neq \emptyset$ دپاره شتون لري.
خو اوس به ئې په ښکاره بيان کړو.

د غړو له مخې د برابرۍ په لرلو سره، که وغواړو $C$ د
$\bigcap\Setabs{X}{\phi(X)}$ په توګه تعريف کړو، په حقيقت کښې
يوازې يو داسې شے $C$ غواړو چې لاندې شرط پوره کړي:
\begin{equation}\ollabel{bicondelimarbintersection}
	\forall x(x \in C \liff \forall X(\phi(X) \lif x \in X))\tag{*}
\end{equation}
اوس فرض کړئ چې \emph{کوم} سټ $S$ شته چې $\phi(S)$ وي.
نو د \olref{bicondelimarbintersection} د ترلاسه کولو دپاره $C$
د \emph{بېلونې} په کارولو سره په ساده توګه داسې تعريفولے شو:
\[
	C = \Setabs{x \in S}{\forall X(\phi(X) \lif x \in X)}.
\]
دا د تمرين په توګه پرېږدو چې وڅېړئ، ايا دا تعريف هماغه غوښتل
شوے \olref{bicondelimarbintersection} ورکوي.
%  fix $x$. First suppose that $x \in C$, as (re)defined; then both $x
%  \in S$ and $\forall X(\phi(X) \lif x \in X)$; but since
%  $\phi(S)$, this reduces to the condition that $\forall X(\phi(X)
%  \lif x \in X)$. Conversely, suppose $\forall X(\phi(X) \lif
%  x \in X)$; then, since $\phi(S)$, we also have that $x \in S$; so
%  $x \in c$, as (re)defined. 
او همدا عمومي لاره په تقاطعګانو تعريفولو کښې د ساده ادراک له هر
ظاهري استعماله ځان ساتل ممکنوي. په هغې ځانګړې بېلګه کښې چې
دا بحث ترې پېل شو، يعنې $\closureofunder{f}{o}$، دا کار داسې
کېږي. د \olref[sfr][infinite][dedekind]{closureproperties} ثبوت مو
په دې يادونه پېل کړے و چې $o \in \ran{f}\cup\{o\}$ او
$\ran{f} \cup \{o\}$ د $f$ له مخې بند دے. نو غوښتل شوے
تعريف داسې ورکولے شو:
\[
	\closureofunder{f}{o} = \Setabs{x \in \ran{f} \cup \{o\}}{(\forall X \ni o)(X \text{ د $f$ له مخې بند دے} \lif x \in X)}.
\]

\end{document}
